\chapter{一段平均操作，为什么可以是一张可逆矩阵}\label{sec:arithmetic}
\section{把参数变化算一次，就知道矩阵从哪里来}
在核心 $K_{p,r}(a,z)$ 中，固定第三块，让前两块作互补混合。第一组拿 $p-k$ 份 $a$ 和 $k$ 份 $rz-a$，所以
\[a'=\frac{p-2k}{p}a+\frac{kr}{p}z,qquad z'=z.
\]
这就是矩阵
\[\binom{a'}{z'}=
\begin{pmatrix}(p-2k)/p&kr/p\\0&1\end{pmatrix}\binom az.
\]
没有增加任何新操作。矩阵只是把两条线性公式写在一起。只要一次或几次平均后重新出现同样的块重数，就能重新读出两个参数，继续用这种记账法。

我们把这样一段有限平均叫作\textbf{返回}，也叫宏。名字里的“返回”是块形状回来，不是数值回到原来。

\section{真正输出与方便的矩阵代表}
回想第5讲的五平均例子，真实输出是 $-(x+y,y)/5$。写成 $U(1)(x,y)$ 更容易看出它是剪切，但实际配置还带着 $-1/5$。

为何可以研究 $U(1)$？因为我们最终要到零。某个后续位置序列能消去 $v$，它也能消去 $cv$，只要 $c\ne0$。因此，在完整配置上，共同尺度不影响可达性。
\begin{definition}\label{def:return}
设 $E_\pi:\Q^2\to\Q^n$ 按分组 $\pi$ 将两个参数展开为等值块配置。若平均词 $W$ 满足
\begin{equation}\label{eq:return}
WE_\pi=\lambda E_{\pi'}M,\qquad \lambda\in\Q^\times,
\end{equation}
且 $\pi'$ 与 $\pi$ 具有相同块重数，称 $M$ 是一个可实现返回。构造须对每个同形状分组适用。可逆矩阵的可实现性按非零共同标量识别。
\end{definition}
把 $E_\pi$ 想成“按照分组名单展开参数”：它将 $(a,z)$ 写到全部 $n$ 个位置。$\pi'$ 是做完后新的分组名单。等式左边先展开再做平均，右边先变参数再按新名单展开，两个结果相同。

\begin{example}[两次正向返回提供逆]
若 $J=\left(\begin{smallmatrix}0&d\\-1&0\end{smallmatrix}\right)$，$d\ne0$，则 $J^2=-dI$。
假如 $J$ 有真实平均词，重复这个词就使全部参数共同乘 $-d$（另计实际词自身的共同尺度）。
所以 $J^{-1}=(-d)^{-1}J$ 在射影意义上由同一条正向词实现。
这不表示一次平均在整个 $\Q^n$ 上可逆；可逆性只发生在带相同块形状的参数空间上。
\end{example}
这不是说某个一般平均投影忽然可逆了。一个任意三元输入平均后丢掉两个差分，当然无法恢复；但我们限制的块族原来就只有两个参数，而且只要求恢复整个配置的比例。额外位置使我们有机会把这两个参数的信息转存到别的块中。

\section{剪切为什么比一般矩阵好用}
记
\[U(t)=\begin{pmatrix}1&t\\0&1\end{pmatrix},qquad
L(t)=\begin{pmatrix}1&0\\t&1\end{pmatrix}.
\]
它们分别给第一坐标加 $t$ 倍第二坐标，或给第二坐标加 $t$ 倍第一坐标。最方便的性质是
\[U(s)U(t)=U(s+t),qquad U(t)^{-1}=U(-t).
\]
如果已经有 $U(2)$ 和 $U(3)$ 以及它们的逆，裴蜀关系 $3-2=1$ 就给 $U(1)$。构造里经常计算若干根参数的最大公因数，原因就在这里。“根子群”在本题中就是让同一类剪切的参数遍历某个加法群。

\section{局部化环不是另一套数，是允许的分母清单}
\begin{definition}[整数的有限局部化]
对整数 $D>1$，
\[
\Z[1/D]=\{a/D^k:a\in\Z,\ k\ge0\}.
\]
等价地，约分后的分母只含 $D$ 的素因子。
其单位是 $\pm\prod_{\ell\mid D}\ell^{e_\ell}$，其中指数 $e_\ell$ 可正可负。
\end{definition}
$\Z[1/6]$ 里有 $1/2,1/3,1/12$，因为它们都能写成 $a/6^k$；没有 $1/5$。数 $5$ 在这个环里，但它的逆不在，所以 $5$ 不是单位。数 $2$ 和 $3$ 则是单位。

在原数值层面，整数输入经 $p$ 平均只能出现 $p$ 的分母。但核心矩阵允许除去一个\emph{整个配置的}共同尺度后表示。它的参数环可能含别的分母。这两句话并不冲突：矩阵中的 $1/2$ 不等于实际坐标被单独免费除以二。

\section{用共轭把一个步长变成另一个步长}
设已取得可逆三角返回 $T$。组合 $TU(t)T^{-1}$ 的意义是：先用 $T^{-1}$ 改参数，再剪切，最后用 $T$ 改回来。计算告诉我们，这相当于换了剪切的步长。
\begin{lemma}[可逆倍率扩展剪切参数]\label{lem:localize}
假设 $U(hR)$ 全部可实现，且上三角矩阵 $T=\left(\begin{smallmatrix}1&b\\0&\kappa\end{smallmatrix}\right)$ 及其逆可实现。则 $U(hR[\kappa,\kappa^{-1}])$ 全部可实现。若 $\kappa=A/B$ 约分，则新环为 $R[1/(AB)]$。
\end{lemma}
\begin{proof}
使用 $TU(t)T^{-1}=U(t/\kappa)$、$T^{-1}U(t)T=U(\kappa t)$ 和 $U(s)U(t)=U(s+t)$。
新环中每个元素是有限个 $r_i\kappa^i$ 的和，其中 $i$ 允许为负数。
先用共轭实现每一项的剪切，再相乘把参数相加，得到所需返回。
全部求和都是有限的，因此最终仍是一条有限平均词。
至于新环中的分母为什么正好来自 $AB$，一方面 $A/B$、$B/A$ 显然都属于 $R[1/(AB)]$；
另一方面，裴蜀等式 $uA+vB=1$ 给出 $u\kappa+v=1/B$；对逆倍率同理得到 $1/A$。
\end{proof}
例如倍率为 $3/2$。它和逆倍率 $2/3$ 允许我们把参数乘这些数。更具体地，$3/2-1=1/2$，$1-2/3=1/3$，所以加法和这些倍率已经生成 $\Z[1/6]$。一般情形中，裴蜀等式承担同一任务。

到此若说“已能实现 $U(hR)$”，意思是对每个属于 $hR$ 的参数，都有某条有限平均词。并非一条词同时实现所有参数，也不只是模几个整数检查到了它们。下一讲将利用这整族剪切，把无界参数之间的精确输送交给群论。
